\Needspace{14\baselineskip}
\section{Experiment 1: Evaluation Results and Low-Uncertainty Subset Extraction}
\label{sec:experiment1_results}
\markboth{\thesection\quad Experiment 1: Evaluation Results}{}

\subsection{Descriptive Statistics}
\label{subsec:descriptive_statistics}

\subsubsection{Accuracy and Hallucination Rate}

The three models produced 6,741 responses in total, of which 6,667 satisfied the XML requirements and entered the formal evaluation, yielding a validity rate of 98.9\%. Cell-level exclusions are reported in Appendix~A.5. This high validity rate indicates that format adaptation reliably supports the subsequent correctness, hallucination, and uncertainty analyses. LLaVA has the lowest overall accuracy (0.493) and the highest hallucination rate (0.679), whereas Qwen and InternVL achieve higher accuracy with substantially lower hallucination rates of 0.358 and 0.322, respectively; InternVL has the strongest overall performance.

\begin{table}[H]
\centering
\caption{Accuracy and hallucination rate in each model--dataset cell.}
\label{tab:accuracy_hallucination_rate}
\datatablefont
\setlength{\tabcolsep}{8pt}
\renewcommand{\arraystretch}{1.16}
\begin{tabular}{llrcc}
\toprule
\tableheadrow
\thead{Model} & \thead{Dataset} & \thead{$n$} & \thead{Accuracy} & \thead{Hallucination rate} \\
\midrule
\multirow{3}{*}{LLaVA} & ViLP & 886 & 0.529 & 0.527 \\
& HallusionBench & 1,108 & 0.514 & \best{0.795} \\
& MM-Vet & 216 & 0.241 & 0.704 \\
\midrule
\tablestriperow \multirow{3}{*}{Qwen} & ViLP & 898 & 0.647 & 0.239 \\
\tablestriperow & HallusionBench & 1,099 & 0.661 & 0.468 \\
\tablestriperow & MM-Vet & 218 & 0.560 & 0.298 \\
\midrule
\multirow{3}{*}{InternVL} & ViLP & 900 & 0.618 & 0.254 \\
& HallusionBench & 1,126 & \best{0.710} & 0.386 \\
& MM-Vet & 216 & 0.551 & 0.264 \\
\midrule
\tableaccentrow LLaVA & \textit{All} & 2,210 & 0.493 & 0.679 \\
\tableaccentrow Qwen & \textit{All} & 2,215 & 0.645 & 0.358 \\
\tableaccentrow InternVL & \textit{All} & 2,242 & \best{0.658} & \best{0.322} \\
\tableaccentrow \textbf{Overall} & & \textbf{6,667} & \textbf{0.599} & \textbf{0.452} \\
\bottomrule
\end{tabular}
\end{table}

HallusionBench places the strongest demands on hallucination control. For example, LLaVA remains above 0.5 in accuracy on this dataset, yet its hallucination rate reaches 0.795. This divergence shows that final-answer capability and factual reliability are related but not interchangeable; accuracy alone does not fully characterize the trustworthiness of an MLLM response.

\subsubsection{Joint Distribution of Correctness and Hallucination}

\begin{figure}[H]
\centering
\includegraphics[width=0.95\linewidth]{fig_3_1_2a_label_joint}
\caption{Joint distribution of correctness and hallucination. The left panel shows the 100\% stacked distribution for each model; the right panel shows hallucination rates conditional on answer correctness.}
\label{fig:label_joint}
\end{figure}

As shown in \Cref{fig:label_joint}, correctness and hallucination are distinct evaluation targets. Among the 6,667 valid responses, 14.0\% are correct but hallucinated and 8.8\% are incorrect without hallucination. Replacing $H$ with the error label $E=1-C$ would therefore mismatch 1,522 samples (22.8\%). Moreover, 30.9\% of all hallucinated responses still have correct final answers, while 22.1\% of all incorrect responses contain no hallucination. Complete counts and conditional proportions are provided in Appendix~A.2. The absolute $\phi$ coefficients in the nine model--dataset cells range only from 0.46 to 0.65, confirming that the labels are moderately to strongly associated but remain far from equivalent.

\begin{figure}[H]
\centering
\includegraphics[width=0.88\linewidth]{fig_3_1_2b_hallucination_types}
\caption{Composition of hallucination types for the three models, shown as three parallel donut charts.}
\label{fig:hallucination_types}
\end{figure}

Across all hallucinated responses, 61.9\% contain both visual and reasoning hallucination, while visual-only and reasoning-only hallucination account for 19.0\% and 19.1\%, respectively (\Cref{fig:hallucination_types}). Hallucination therefore often propagates through the visual-description--reasoning--answer chain rather than appearing only in the final answer, which further motivates independent hallucination labels.

\subsection{Error and Hallucination Detection Performance}
\label{subsec:detection_performance}

\subsubsection{Error Detection}

\begin{figure}[H]
\centering
\includegraphics[width=0.96\linewidth]{fig_3_2_1_error_detection}
\caption{Error-detection AUROC of the three UQ methods across models and datasets. Error bars show 95\% bootstrap confidence intervals; the dashed line marks random performance at 0.5.}
\label{fig:error_detection}
\end{figure}

As shown in \Cref{fig:error_detection}, Semantic Entropy achieves the best AUROC in eight of the nine model--dataset cells and has a nine-cell mean of 0.724. UMPIRE is best only for LLaVA/HallusionBench, while PPL is not best in any cell. The mean ranks of SE, UMPIRE, and PPL are 1.22, 2.11, and 2.67, respectively. The model-wise macro-average SE AUROC is also consistently highest (0.672, 0.729, and 0.771 for LLaVA, Qwen, and InternVL), while all three methods generally perform better on MM-Vet than on ViLP and HallusionBench. Full cell-level AUROC values and confidence intervals are reported in Appendix~A.3.

\subsubsection{Hallucination Detection}

\begin{figure}[H]
\centering
\includegraphics[width=0.96\linewidth]{fig_3_2_2_hallucination_detection}
\caption{Hallucination-detection AUROC of the three UQ methods across models and datasets. Error bars show 95\% bootstrap confidence intervals; the dashed line marks random performance at 0.5.}
\label{fig:hallucination_detection}
\end{figure}

Changing the target from error to hallucination weakens all three methods (\Cref{fig:hallucination_detection}). The nine-cell mean AUROC decreases from 0.658, 0.724, and 0.685 for error detection to 0.568, 0.655, and 0.613 for PPL, SE, and UMPIRE, respectively. SE remains best in eight cells, but its absolute performance is lower than for error detection. Four PPL confidence intervals and one UMPIRE confidence interval cover random performance at 0.5, whereas none of the nine SE intervals does. The model-wise macro-average SE AUROC is 0.623, 0.636, and 0.707 for LLaVA, Qwen, and InternVL, all below the corresponding error-detection results. These results indicate that semantic dispersion across samples is more effective at capturing answer instability than factual hallucination in visual observation and reasoning. Full results are reported in Appendix~A.4.

\subsubsection{Detection-Target Gap}

\begin{table}[H]
\centering
\caption{Mean $\operatorname{AUROC}(H)-\operatorname{AUROC}(E)$ across the three datasets for each model and UQ method.}
\label{tab:target_gap}
\datatablefont
\setlength{\tabcolsep}{12pt}
\renewcommand{\arraystretch}{1.16}
\begin{tabular}{lccc}
\toprule
\tableheadrow
\thead{Model} & \thead{PPL} & \thead{SE} & \thead{UMPIRE} \\
\midrule
LLaVA & $-0.061$ & $-0.049$ & $-0.045$ \\
\tablestriperow Qwen & \best{$-0.126$} & \best{$-0.093$} & \best{$-0.105$} \\
InternVL & $-0.083$ & $-0.063$ & $-0.066$ \\
\bottomrule
\end{tabular}
\end{table}

All nine model--method means in \Cref{tab:target_gap} are negative, with an overall mean of $-0.077$. Qwen exhibits the largest decrease for all three methods, whereas LLaVA has the smallest. Complete cell-level paired differences and 95\% bootstrap confidence intervals are provided in Appendix~A.4.1. When the analysis is restricted to incorrect answers and asks whether an incorrect response is also hallucinated, the nine-cell mean AUROC is only 0.400 for PPL, 0.494 for SE, and 0.423 for UMPIRE; 22 of 27 combinations lie below 0.5, with a minimum of 0.307 for UMPIRE on LLaVA/HallusionBench. Independent hallucination labels therefore cannot be replaced by answer-error labels even within the error subset.

\subsection{Low-Uncertainty Hallucination Blind Spot}
\label{subsec:luh_blind_spot}

For each model and UQ method, I pool the three datasets and use the $\alpha$-quantile of the non-hallucinated score distribution as the threshold:
\[
\operatorname{LUH\_share}(\alpha)
=
\Pr\!\left(s_{H=1}\le Q_{\alpha}(s_{H=0})\right),
\qquad \alpha\in\{0.25,0.50\}.
\]
The threshold is determined only from non-hallucinated samples, and all ties at the threshold are included. The value $\alpha=0.25$ defines the stricter lower-quartile blind spot, whereas $\alpha=0.50$ defines the broader median blind spot.

\subsubsection{Blind-Spot Share and Missed-Hallucination Count}

\begin{table}[H]
\centering
\caption{LUH share and absolute count at the two quantile thresholds.}
\label{tab:luh_share_count}
\datatablefont
\setlength{\tabcolsep}{5.2pt}
\renewcommand{\arraystretch}{1.15}
\begin{tabular}{lr l rr rr}
\toprule
\tableheadrow
\multirow{2}{*}{\thead{Model}} & \multirow{2}{*}{\thead{$n_H$}} & \multirow{2}{*}{\thead{Method}} & \multicolumn{2}{c}{\thead{$\alpha=0.25$}} & \multicolumn{2}{c}{\thead{$\alpha=0.50$}} \\
\cmidrule(lr){4-5}\cmidrule(l){6-7}
\tableheadrow
& & & \thead{Share} & \thead{Count} & \thead{Share} & \thead{Count} \\
\midrule
\multirow{3}{*}{LLaVA} & \multirow{3}{*}{1,500} & PPL & 0.195 & 292 & 0.434 & 651 \\
& & SE & \best{0.137} & \best{206} & 0.453 & 680 \\
& & UMPIRE & 0.169 & 254 & \best{0.397} & \best{595} \\
\midrule
\tablestriperow \multirow{3}{*}{Qwen} & \multirow{3}{*}{794} & PPL & 0.246 & 195 & 0.515 & 409 \\
\tablestriperow & & SE & 0.236 & 187 & \best{0.301} & \best{239} \\
\tablestriperow & & UMPIRE & \best{0.164} & \best{130} & 0.385 & 306 \\
\midrule
\multirow{3}{*}{InternVL} & \multirow{3}{*}{721} & PPL & 0.166 & 120 & 0.449 & 324 \\
& & SE & 0.158 & 114 & \best{0.268} & \best{193} \\
& & UMPIRE & \best{0.062} & \best{45} & 0.312 & 225 \\
\bottomrule
\end{tabular}
\end{table}

At $\alpha=0.25$, LUH shares range from 0.062 to 0.246, corresponding to 45--292 missed hallucinations. At $\alpha=0.50$, the shares increase to 0.268--0.515 and the counts to 193--680. Except for Qwen/PPL at the median threshold, the shares remain below their corresponding $\alpha$, indicating useful ranking ability but a substantial residual blind spot. Each row defines a method-specific LUH subset; subsets may overlap and must not be added across methods.

\subsubsection{Failure Attribution}

The three UQ methods exhibit distinct structural failure patterns:

\begin{itemize}[leftmargin=2em]
\item \textbf{PPL: answer form and length.} Missed LUH responses are often short Yes/No answers, numbers, or phrases of at most three words, with PPL close to 1.000. Within hallucinated samples, the Spearman correlation between answer-token count and PPL is 0.122 for LLaVA, 0.187 for Qwen, and 0.389 for InternVL. Length normalization and the probability floor of short answers can therefore make an incorrect answer appear highly certain.
\item \textbf{SE: semantic-cluster collapse.} Missed LUH samples typically contain one dominant semantic cluster with mass close to 1 and SE close to 0. SE detects unstable meanings but not the stable repetition of the same hallucination.
\item \textbf{UMPIRE: jointly self-consistent representation and probability.} Missed LUH samples have smaller semantic volume and near-zero mean incoherence. The internal representation and generation probability support the same answer without an external-evidence constraint capable of correcting it.
\end{itemize}

Low-uncertainty hallucination is therefore not a single random failure but a systematic interaction among high-probability short answers, stable semantic clusters, and internally consistent representations. This motivates matching positives and negatives on all three UQ scores before evaluating whether a new method uses information beyond the baselines.

\subsection{Representative Case Analysis}
\label{subsec:representative_cases}

The following cases illustrate three mechanisms: a language prior overriding visual evidence, a visual readout error, and a correct final answer accompanied by a hallucinated intermediate observation.

\begin{table}[!tbp]
\centering
\refstepcounter{table}\label{tab:luh_case1}
\begin{casecard}
\begin{minipage}[t]{0.70\linewidth}
\textbf{Case \thetable(a): Language prior overrides visual evidence}\par\smallskip
{\comicfont\textbf{Q:} The kangaroo is known for its excellent jumping. Which animal native to Australia is shown in the image demonstrating a jump?}
\end{minipage}\hfill
\begin{minipage}[t]{0.25\linewidth}
\raggedleft\small
\textbf{Model:} LLaVA\\
\textbf{Dataset:} ViLP\\
\textbf{Ground truth:} Koala
\end{minipage}

\medskip
\begin{minipage}[t]{0.34\linewidth}
\vspace{0pt}\centering
\includegraphics[width=\linewidth]{vilp_55_case2}
\par\smallskip
{\footnotesize ViLP row 55, case 2}
\end{minipage}\hfill
\begin{minipage}[t]{0.62\linewidth}
\vspace{0pt}\footnotesize\raggedright\comicfont
\textbf{\texttt{<vision>}} The image shows a large animal with a \colorbox{yellow!28}{\strut long tail}, large paws, and a distinctive coat pattern leaping into the air. Its body is horizontal, with extended legs, a raised head, a distinctive muzzle, and large eyes. \textbf{\texttt{</vision>}}

\medskip
\textbf{\texttt{<reasoning>}} The animal in the image is a \colorbox{orange!22}{\strut kangaroo}, a native Australian mammal known for jumping. \textbf{\texttt{</reasoning>}}

\medskip
\textbf{\texttt{<answer>}} kangaroo \textbf{\texttt{</answer>}}

\medskip
\textbf{Judge:} Rating 1; \texttt{vision\_hallucination} + \texttt{reasoning\_hallucination}. The response mistakes a koala-like animal for a kangaroo and invents a long tail.
\end{minipage}

\medskip
\textbf{Analysis.} The question's prior about kangaroos overrides the image evidence, causing both the visual description and reasoning to introduce unsupported features.
\end{casecard}
\end{table}

\begin{table}[!tbp]
\centering
\phantomsection\label{tab:luh_case2}
\begin{casecard}
\begin{minipage}[t]{0.70\linewidth}
\textbf{Case \thetable(b): Visual readout error}\par\smallskip
{\comicfont\textbf{Q:} What is the gasoline price?}
\end{minipage}\hfill
\begin{minipage}[t]{0.25\linewidth}
\raggedleft\small
\textbf{Model:} Qwen\\
\textbf{Dataset:} MM-Vet\\
\textbf{Ground truth:} 4.299
\end{minipage}

\medskip
\begin{minipage}[t]{0.34\linewidth}
\vspace{0pt}\centering
\includegraphics[width=\linewidth]{mmvet_v1_16}
\par\smallskip
{\footnotesize MM-Vet \texttt{v1\_16}; OCR and mathematics}
\end{minipage}\hfill
\begin{minipage}[t]{0.62\linewidth}
\vspace{0pt}\footnotesize\raggedright\comicfont
\textbf{\texttt{<vision>}} The gas-pump display shows ``This Sale.'' Below it, 62.85 is presented as the price, and the gallons dispensed are read as 4.620. \textbf{\texttt{</vision>}}

\medskip
\textbf{\texttt{<reasoning>}} The gasoline price is explicitly displayed as \colorbox{orange!22}{\strut 62.85}. \textbf{\texttt{</reasoning>}}

\medskip
\textbf{\texttt{<answer>}} 62.85 \textbf{\texttt{</answer>}}

\medskip
\textbf{Judge:} Rating 1; \texttt{vision\_hallucination} + \texttt{reasoning\_hallucination}. The value 62.85 is the total sale, the unit price is approximately 4.299 dollars per gallon, and the gallons display is 14.620 rather than 4.620.
\end{minipage}

\medskip
\textbf{Analysis.} A local OCR/readout error propagates through the reasoning chain and produces a confidently worded but visually unsupported answer.
\end{casecard}
\end{table}

\begin{table}[!tbp]
\centering
\phantomsection\label{tab:luh_case3}
\begin{casecard}
\begin{minipage}[t]{0.70\linewidth}
\textbf{Case \thetable(c): Correct answer with a hallucinated observation}\par\smallskip
{\comicfont\textbf{Q:} From 2012 to 2022, did Apple derive the majority of its revenue from Japan?}
\end{minipage}\hfill
\begin{minipage}[t]{0.25\linewidth}
\raggedleft\small
\textbf{Model:} InternVL\\
\textbf{Dataset:} HallusionBench\\
\textbf{Ground truth:} No
\end{minipage}

\medskip
\begin{minipage}[t]{0.38\linewidth}
\vspace{0pt}\centering
\includegraphics[width=\linewidth]{hb_chart_6_1_3}
\par\smallskip
{\footnotesize VS/chart, set 6, figure 1, question 3}
\end{minipage}\hfill
\begin{minipage}[t]{0.58\linewidth}
\vspace{0pt}\footnotesize\raggedright\comicfont
\textbf{\texttt{<vision>}} The chart shows Apple's revenue share by region from 2012 to 2022. The Japan region is \colorbox{yellow!28}{\strut consistently the smallest} throughout the entire period. \textbf{\texttt{</vision>}}

\medskip
\textbf{\texttt{<reasoning>}} Since Japan is the smallest and never exceeds 10\% of total revenue, Apple did not derive the majority of its revenue from Japan. \textbf{\texttt{</reasoning>}}

\medskip
\textbf{\texttt{<answer>}} no \textbf{\texttt{</answer>}}

\medskip
\textbf{Judge:} Rating 2; \texttt{vision\_hallucination}; \textbf{correct=true}. Rest of Asia Pacific is often smaller than Japan, so the highlighted absolute claim is inaccurate.
\end{minipage}

\medskip
\textbf{Analysis.} The final answer is correct, but the intermediate visual observation is inaccurate. Correctness and hallucination must therefore be labeled and evaluated separately.
\end{casecard}
\end{table}

\FloatBarrier

Together, the cases show that checking only the final answer is insufficient for determining whether the whole response is supported by visual evidence. The attribution analysis above explains why the baselines can assign low uncertainty: high token probability for short or numeric answers lowers PPL, stable repetition of an erroneous conclusion lowers SE, and agreement between internal representations and generation probabilities prevents UMPIRE from identifying content that lacks visual support.

\subsection{Low-Uncertainty Difficult Subset}
\label{subsec:luh_subset}

\subsubsection{Construction Procedure}

The difficult subset constructs positive and negative samples that receive similar scores from the three UQ methods but have opposite hallucination labels:

\begin{enumerate}[label=\textbf{Step \arabic*:}, leftmargin=5.2em, labelsep=0.7em]
\item For each model, build a candidate pool containing valid judge labels, valid PPL/SE/UMPIRE scores, a non-missing hallucination label, and an image input.
\item Within every model--dataset cell, convert the three UQ scores to percentiles and compute $\bar p=\frac{1}{3}(p_{\mathrm{PPL}}+p_{\mathrm{SE}}+p_{\mathrm{UMPIRE}})$.
\item Pool the three datasets within each model and select the 200 hallucinated samples with the lowest $\bar p$ as LUH positives.
\item Using the three score percentiles as matching features, greedily match each positive without replacement to the nearest non-hallucinated sample in the same model by Euclidean distance, processing positives in ascending $\bar p$ order.
\item Verify labels, image inputs, and sample counts, yielding 200 LUH positives and 200 matched non-hallucinated negatives per model.
\end{enumerate}

Candidate-pool sizes, score alignment, matching distances, and final subset composition are reported in Appendix~A.4.2.

\subsubsection{Baseline Performance}

\begin{table}[H]
\centering
\caption{Baseline performance on the difficult subsets.}
\label{tab:subset_baselines}
\datatablefont
\setlength{\tabcolsep}{6pt}
\renewcommand{\arraystretch}{1.16}
\begin{tabular}{lcccccc}
\toprule
\tableheadrow
\multirow{2}{*}{\thead{Model}} & \multicolumn{2}{c}{\thead{PPL}} & \multicolumn{2}{c}{\thead{SE}} & \multicolumn{2}{c}{\thead{UMPIRE}} \\
\tableheadrow
& \thead{AUROC} & \thead{AUPRC} & \thead{AUROC} & \thead{AUPRC} & \thead{AUROC} & \thead{AUPRC} \\
\cmidrule(lr){2-3}\cmidrule(lr){4-5}\cmidrule(l){6-7}
LLaVA & \best{0.402} & 0.435 & 0.476 & 0.464 & \best{0.413} & 0.437 \\
\midrule
\tablestriperow Qwen & 0.485 & 0.499 & 0.499 & 0.496 & 0.477 & 0.489 \\
InternVL & 0.489 & 0.498 & 0.498 & 0.500 & \best{0.502} & \best{0.502} \\
\bottomrule
\end{tabular}
\end{table}

Across the nine model--method combinations, AUROC ranges from 0.402 to 0.502 with a mean of 0.471, and eight of nine values are below 0.5. For LLaVA, the upper confidence bounds of PPL and UMPIRE are both below 0.5, indicating stable mild reverse ordering; the Qwen and InternVL scores are near random and most confidence intervals cover 0.5. After aligning the score distributions of all three UQ methods, PPL, SE, and UMPIRE cannot reliably distinguish LUH from non-hallucinated responses. The subset therefore provides a strict baseline for evaluating ERA in Experiment~2.

In summary, Experiment~1 provides a progression of evidence from label relationships to detection performance and finally to failure cases. Correctness and hallucination are correlated but not interchangeable. PPL, SE, and UMPIRE rank answer errors to some extent, yet perform consistently worse on hallucination and miss a nontrivial set of low-uncertainty hallucinations. After their score distributions are explicitly aligned in the difficult subsets, all three methods approach random AUROC, showing that generation probability, semantic consistency across samples, and internally self-consistent representations do not guarantee sufficient support from the image and question. Experiment~2 therefore shifts the focus from how certain the model is to whether that certainty is grounded in external evidence, motivating early cross-modal attribution signals.
